% =============================================================================
% chapter0_intro.tex — General Introduction
% =============================================================================

\chapter*{General Introduction}
\addcontentsline{toc}{chapter}{General Introduction}
\markboth{General Introduction}{General Introduction}
\thispagestyle{fancy}

The data analytics industry is undergoing a profound and irreversible shift.
For three decades, the Statistical Analysis System (SAS) dominated enterprise
data processing: its procedural execution model, built-in statistical procedures,
and tightly integrated macro language made it the tool of choice for financial
institutions, pharmaceutical companies, and government agencies. Today, that
dominance is eroding. Open-source Python --- with its pandas, NumPy, and PySpark
ecosystem --- offers comparable analytical power at a fraction of the licensing
cost, with a far larger developer community and native integration into
cloud-native data platforms.

The consequence is a migration wave. Organisations are faced with the task of
converting years of accumulated SAS code into Python, and the scale of this
task is daunting. Manual rewriting is slow, error-prone, and prohibitively
expensive at enterprise scale. Existing automation tools partially address the
problem: rule-based transpilers handle simple constructs but fail on the macro
meta-programming that permeates enterprise SAS codebases; raw Large Language
Model (LLM) translation produces plausible-looking Python that silently diverges
from the SAS semantics it was meant to replicate.

This project presents \codara{} --- a Multi-Agent SAS-to-Python/PySpark
Conversion Accelerator. \codara{} is built on the observation that the
migration problem has not one but five distinct sub-problems, each requiring a
different technical approach: (1) file-level parsing and cross-file dependency
resolution; (2) semantic partitioning of SAS code into translatable units;
(3) complexity-aware retrieval of relevant translation examples; (4) LLM-driven
translation with structured output; and (5) formal and behavioural verification
that the translated code is semantically equivalent to the original.

The system decomposes these five concerns across an eight-node LangGraph
orchestration pipeline, three complementary Retrieval-Augmented Generation
(RAG) paradigms, Z3 Satisfiability Modulo Theories (SMT) formal verification,
the Constraint-Driven Adversarial Input Synthesis (CDAIS) framework, and the
SemantiCheck four-layer semantic scoring system. The result is a
production-ready system deployed as a containerised service on Azure Container
Apps, with a React web dashboard, a REST API, and a six-job CI/CD pipeline.

\textbf{Report organisation.} This report is structured as follows:

\begin{itemize}
  \item \textbf{Chapter~\ref{ch:context} --- Project Context and Foundation}
    situates the work within its academic and industrial frames, characterises
    the SAS migration problem, surveys existing tools and their limitations,
    and justifies the hybrid DSR + Agile methodology.

  \item \textbf{Chapter~\ref{ch:foundations} --- Theoretical Foundations and
    Related Work} reviews the scientific principles underlying every major
    system component: SAS language semantics, automated code translation,
    three RAG paradigms, multi-agent orchestration, SMT formal verification,
    program slicing, adversarial test synthesis, composite verification scoring,
    and ML calibration.

  \item \textbf{Chapter~\ref{ch:specifications} --- Project Specifications}
    formalises the business, technical, and scientific objectives; derives
    functional and non-functional requirements; and presents use cases and
    a high-level architecture.

  \item \textbf{Chapter~\ref{ch:architecture} --- System Architecture and
    Design Decisions} documents the empirical technology-selection process,
    the eight-node pipeline design, the multi-store data architecture, and
    the gold-standard corpus construction.

  \item \textbf{Chapter~\ref{ch:implementation} --- Knowledge Integration and
    Pipeline Implementation} details every component of the implemented system:
    the three-tier RAG strategy, RAPTOR clustering, complexity scoring,
    translation and validation, Z3 verification, CDAIS adversarial synthesis,
    SemantiCheck, and scenario testing.

  \item \textbf{Chapter~\ref{ch:deployment} --- Deployment and User Interface}
    covers the FastAPI backend, service-layer architecture, Docker containerisation,
    Azure Key Vault integration, GitHub Actions CI/CD, and the React frontend.

  \item \textbf{Chapter~\ref{ch:evaluation} --- Evaluation, Results, and
    Discussion} presents all quantitative results against the project targets,
    analyses the system's strengths and limitations, situates the contributions
    within Design Science Research, and proposes future research directions.
\end{itemize}
